\problem{}
\par \textbf{1-IN-3 SAT:} Let 1-IN-3 SAT denote the following decision problem: given a Boolean formula $\phi$ in 3CNF, decide whether $\phi$ has a satisfying assignment such that every clause contains exactly one true literal. Prove that 1-IN-3 SAT is NP-complete.

\par \textbf{Hint:} Reduce from 3-SAT. Create for every clause 8 new variables $y_{000}, y_{001}$, $\dots$, $y_{111}$ and reduce 3-SAT to 1-IN-AT-MOST-7 SAT. The idea is that precisely 1 of these new variables should be true for each assignment of the clause.

\noindent \textbf{Proof:}

\noindent\textbf{Step 1: \(1-IN-3 SAT \in \mathrm{NP}\).} \\
A certificate is a Boolean assignment to all variables.
A verifier checks that in every clause of $\varphi$ exactly one literal
is true. This requires constant time per clause and therefore runs in
polynomial time. Hence the problem is in NP.


\noindent\textbf{Step 2: Reduction from \textsc{Partition} to \(1-IN-3 SAT\).} \\

We reduce from the NP-complete problem \textsc{3-SAT}.
Let the input formula be
\[
\varphi = C_1 \land C_2 \land \cdots \land C_m,
\]
where each clause $C_j = (\ell_{j1} \lor \ell_{j2} \lor \ell_{j3})$
contains three literals.

For each clause $C_j$ we introduce eight new Boolean variables
\[
y_{000}^{(j)},\; y_{001}^{(j)},\; y_{010}^{(j)},\; \dots,\;
y_{111}^{(j)},
\]
indexed by all triples $(a,b,c) \in \{0,1\}^3$, which represent the
possible truth-value patterns of the three literals in $C_j$.  Intuitively,
$y_{abc}^{(j)} = 1$ is intended to mean that under the chosen assignment
the three literals $(\ell_{j1}, \ell_{j2}, \ell_{j3})$ evaluate
respectively to $(a,b,c)$.

We enforce the following constraints using a polynomial number of 
1-in-3 clauses:

\begin{itemize}
    \item \emph{Exactly one pattern is chosen for each clause.}
    For the eight variables $y_{000}^{(j)}, \dots, y_{111}^{(j)}$ we add
    a collection of 1-in-3 clauses that together enforce that
    \emph{exactly one} of them is true. This uses the fact that
    ``exactly-one over $t$ variables'' can be encoded with $O(t)$
    1-in-3 clauses and auxiliary variables.

    \item \emph{Consistency with original literals.}
    For each $y_{abc}^{(j)}$ we enforce that if it is true,
    then $\ell_{j1}$ has value $a$, $\ell_{j2}$ has value $b$, and
    $\ell_{j3}$ has value $c$. Each implication
    $y_{abc}^{(j)} =1 \Rightarrow (\ell_{jk} = a)$
    is encoded with a constant number of 1-in-3 clauses using auxiliary
    variables. Thus all consistency relations are expressible in
    polynomially many 1-in-3 clauses.
\end{itemize}

To ensure that each chosen pattern corresponds to a clause being
satisfied in the original formula (i.e., that at least one literal is
true), we simply omit the pattern $y_{000}^{(j)}$ from the construction
(or equivalently forbid it with clauses).
Thus for each clause we keep only seven pattern variables
$y_{001}, y_{010}, \dots, y_{111}$,
and enforce that exactly one of these seven is true.
This guarantees that in the chosen triple at least one literal is true.

The total number of new variables and clauses is $O(m)$ and the entire
construction runs in polynomial time.


\noindent\textbf{Step 3: Correctness (equivalence).} \\

\noindent ($\Rightarrow$)
Suppose $\varphi$ is satisfiable. Let $x^*$ be a satisfying assignment.
For each clause $C_j$ let $(a,b,c)$ be the truth values of its literals
$(\ell_{j1},\ell_{j2},\ell_{j3})$ under $x^*$.
Since $C_j$ is satisfied, at least one of $a,b,c$ equals 1.
We then set $y_{abc}^{(j)} = 1$ and all other pattern variables for
clause $j$ to 0.  
All consistency constraints are satisfied.
Thus the constructed 1-IN-3 instance is satisfied by setting exactly one 
pattern variable per clause to true, proving the instance is a 
yes-instance of 1-IN-3 SAT.


\noindent ($\Leftarrow$)
Conversely, suppose the constructed 1-IN-3 SAT instance is satisfiable.
For each clause $C_j$, exactly one pattern variable $y^{(j)}_{abc}$ is
true. Because we forbid $y_{000}$, we must have $(a,b,c)\neq(0,0,0)$.
The consistency clauses enforce that the literals of $C_j$ take truth
values $(a,b,c)$ in the induced assignment to the original variables.
Thus at least one literal in each clause is true, and therefore the
original 3-SAT formula $\varphi$ is satisfied.

Therefore 1-IN-3 SAT is NP-complete.

\newpage